%% pc-reflexion-transmission — répartition de l'énergie d'une onde à la traversée d'une paroi.
%% Onde incidente I (solideA) -> réfléchie I_r (solideB), absorbée I_a (solideD), transmise I_t (solideE).
%% Le motif de briques est dessiné à la main (pas de `pattern`, trop lourd une fois converti en SVG).
\documentclass[tikz,border=4pt]{standalone}
\usepackage{liaisons}
\usetikzlibrary{decorations.pathmorphing}
\begin{document}
\begin{tikzpicture}[x=1cm,y=1cm,
  onde/.style={-{Stealth[length=6pt]}, line width=1.3pt},
  txt/.style={font=\scriptsize},
  etiq/.style={font=\scriptsize, align=center, inner sep=1.5pt}]

%% ---------------- la paroi (milieu B), motif de briques --------------------
\def\pl{0}\def\pr{2.0}\def\pb{-1.5}\def\pt{2.0}
\begin{scope}
  \path[clip] (\pl,\pb) rectangle (\pr,\pt);
  \fill[solideD!15] (\pl,\pb) rectangle (\pr,\pt);
  \foreach \k in {0,1,...,10} {
    \draw[black!45, line width=0.4pt] (\pl,{\pb+0.35*\k}) -- (\pr,{\pb+0.35*\k});
    \foreach \j in {0,1,2,3} {
      \draw[black!45, line width=0.4pt]
        ({\pl+0.7*\j+0.35*mod(\k,2)},{\pb+0.35*\k}) -- ++(0,0.35); } }
\end{scope}
\draw[line width=0.8pt] (\pl,\pb) rectangle (\pr,\pt);
\node[txt, anchor=south] at (1.0,\pt+0.08) {\textbf{Milieu B}};
\node[txt, anchor=center] at (-1.55,0.78) {\textbf{Milieu A}};

%% ---------------- onde incidente et onde réfléchie -------------------------
\coordinate (P) at (0,0.4);                 % point d'impact sur la face gauche
\draw[onde, solideA] (-2.2,-0.9) -- (P);
\draw[onde, solideB] (P) -- (-2.2,1.7);
\draw[black!45, line width=0.4pt, dash pattern=on 2pt off 2pt] (-1.05,0.4) -- (0.5,0.4);
\node[etiq, anchor=north west, text=solideA] at (-2.25,-0.95) {onde incidente\\$I$};
\node[etiq, anchor=south west, text=solideB] at (-2.25,1.72) {\textbf{Réflexion} $I_r$};

%% ---------------- ondes transmises -----------------------------------------
\foreach \yy in {1.15,0.4,-0.35} {
  \draw[onde, solideE] (\pr,0.4) -- (3.35,\yy); }
\node[etiq, anchor=west, text=solideE] at (3.45,0.4) {\textbf{Transmission}\\$I_t$};

%% ---------------- absorption dans la paroi ---------------------------------
\foreach \xx in {0.45,1.0,1.55} {
  \draw[solideD, line width=0.9pt, -{Stealth[length=4pt]},
        decorate, decoration={snake, amplitude=0.7pt, segment length=4pt, post length=3pt}]
       (\xx,0.15) -- (\xx,-0.85); }
\fill[solideD] (\pl,\pb) rectangle (\pr,-1.0);
\node[txt, text=white, anchor=center] at (1.0,-1.25) {\textbf{Absorption} $I_a$};

%% ---------------- relation encadrée ----------------------------------------
\rappel{\node[draw, line width=0.7pt, rounded corners=2pt, inner sep=4pt, align=center,
      anchor=north] at (1.3,-1.85)
  {$I = I_r + I_a + I_t$\\[1pt]
   {\scriptsize l'énergie incidente se répartit intégralement entre les trois parts}};}
\end{tikzpicture}
\end{document}
